\scanpage{102}
la valeur \unsure{correcte|voulue|correct} pour \deleted{\unsure{le|les}} $G/G^0$ : \deleted{\unsure{G}} en effet, il est dual au groupe des racines de l'unité, or ce dernier est bien le dual de $\widehat{\mathbb Z}$ !

\margin{$M(q^m)\subset M(\underset{\scriptstyle \Vert\atop\scriptstyle p^n}{q})\subset M(p)$
\par $M(q^m)^m\subset M(q)^m\subset M(q^{\deleted{\unsure{m}}m})\subset M(q)$}

Soit $M(q)$ le ss-groupe de $Q'(i)^*$ qu'on définit en termes de \unsure{$q$} ci-dessus, soit \deleted{\unsure{$f$}} $f_q\in G(q)$ son dual, qui est le groupe de Galois motivique de $\mathbb F_q$. Considérons l'extension $\mathbb F_{q^m}$ de $\mathbb F_q$, d'où

% slanted
\margin{$M(q)\simeq\mathbb Z\times M^0\deleted{(q)}$ si $q=p^n$, $n$ pair. $M^0\deleted{\unsure{(q)}}$ formé des $u\in Q'(i)^*$ qui sont des unités \unsure{abs.} exc.\,$p$, et vérifient \deleted{\unsure{à}}\deleted{\unsure{$*$}} $v_{\mathfrak p}(u)+v_{\mathfrak p^{\unsure{c|\prime}}}(u)=0$ ; le \deleted{quotient} \deleted{\unsure{de}} \deleted{\unsure{$M$}} $\mathbb Z$ correspond à $q^{\unsure{n|1}/2}$. \deleted{\unsure{les}} les applications de transition entre les $M(q)$ sont l'identité sur le \deleted{facteur} $\mathbb Z$, \deleted{défini} et la multiplication par $m$ sur le facteur $M^0$.}

\deleted{\unsure{$G(q^m$|$G/$}}
\diagram{\begin{tikzcd}[ampersand replacement=\&]
G(q) \& G(q^m)\arrow[l]\\
(f_q)^m \& f_{q^m}\arrow[l,decorate,decoration={snake,amplitude=.35mm,segment length=2.5mm,post length=.9mm}]
\end{tikzcd}}
;

\deleted{\unsure{On}}\ \deleted{\unsure{définit}} on voit que l'hom. transposé n'est autre que
\diagram{\begin{tikzcd}[ampersand replacement=\&]
M(q)\arrow[r,"\text{élévation à la puissance }m\text{ième}"] \& M(q^m)
\end{tikzcd}}
dont le noyau est (comme de juste) \deleted{\unsure{$\mu$}} $\mu_m$ (dual de $\mathbb Z/m\mathbb Z$) et dont l'image est $M(q)^m$ donc contient $M(q^m)^m$, dont le conoyau est annulé par $m$. Par suite, le \deleted{conoyau} noyau de \unsure{$\varphi$} est annulé par $m$.

\deleted{Consi} Considérons la limite inductive des $M(q)$ ($q$ variant multiplicativement) par les homs précédents : on trouve $M(p^\infty)$, un groupe sans torsion dont le dual doit être le groupe de
% archive sheet number: 101

